\graphicspath{{lectures/07_eloadas_abrak/}}

\section{7. előadás -- Lineáris differenciálegyenlet-rendszerek és magasabbrendű egyenletek}

\subsection{Lineáris differenciálegyenlet-rendszerek (folytatás)}
\textit{EA: 2019. 11. 06. (7). A kézirat hivatkozása: „L. An. alk. 7. ea.”}

\textbf{Emlékeztető.} Lineáris differenciálegyenlet-rendszer:
\[
 x'(t)=A(t)x(t)+b(t),\qquad t\in I.
\]
\textit{(Volt: kezdetiérték-probléma; a kezdetiérték-probléma megoldásának létezése és egyértelműsége.)}

\subsection{A homogén lineáris rendszer megoldástere}
\textbf{Homogén egyenlet:} $b(t)=0$, tehát $x'=Ax$.

\begin{matbox}
\textbf{Tétel.} Jelöljük az összes teljes homogén megoldás halmazát:
\[
 \mathcal M_h:=\{\varphi:I\to\mathbb R^n\mid
 \varphi'=A\varphi,\ \varphi\text{ teljes megoldás}\}.
\]
\begin{enumerate}[label=\arabic*)]
\item $\mathcal M_h\subset C^1(I,\mathbb R^n)$ egy $n$ dimenziós altér.
\item Ha $\varphi_1,\varphi_2,\ldots,\varphi_n\in\mathcal M_h$, akkor
\[
 \begin{gathered}
 \varphi_1,\ldots,\varphi_n\ \text{lineárisan függetlenek}\\
 \Longleftrightarrow\\
 \exists t_0\in I:\quad \varphi_1(t_0),\ldots,\varphi_n(t_0)\in\mathbb R^n\\
 \text{lineárisan függetlenek.}
 \end{gathered}
\]
\end{enumerate}
\end{matbox}

\textbf{Bizonyítás.} \textit{(Vázlat; a kézirat hangsúlyozza: „Ez kell!!!”)}
\begin{enumerate}[label=\arabic*)]
\item $\mathcal M_h$ altér (triviális).
\item $\dim\mathcal M_h=n$. Belátjuk a következő kezdetiérték-problémákból:
\[
 x'=Ax,\qquad x(\tau)=e_i,\qquad i=1,\ldots,n,
\]
ahol az $\mathbb R^n$ kanonikus bázisának vektorai:
\[
 e_i=(0,\ldots,0,\underbrace{1}_{i\text{-edik hely}},0,\ldots,0)^T.
\]
\end{enumerate}
\textit{Kéziratos hivatkozás: 07\_An\_alk\_2022\_osz.pdf.}

\subsubsection*{Alaprendszer és alapmátrix}
\textbf{Definíció.}
\begin{enumerate}[label=\arabic*)]
\item Ha $\mathcal M_h$ egy $\varphi_1,\ldots,\varphi_n$ bázisát
vesszük az $x'=Ax$ homogén egyenletben, akkor ez a bázis
\textbf{alaprendszer}.
\item Az alaprendszer oszlopaiból álló mátrix:
\[
 \Phi(t):=\bigl(\varphi_1(t),\ldots,\varphi_n(t)\bigr)
 =\begin{bmatrix}\vert&&\vert\\
                   \varphi_1(t)&\cdots&\varphi_n(t)\\
                   \vert&&\vert\end{bmatrix}
 \in\mathbb R^{n\times n},
\]
\[
 \Phi:I\longrightarrow\mathbb R^{n\times n}
 \qquad\textbf{alapmátrix}.
\]
\end{enumerate}

\begin{matbox}
\textbf{Tétel.} Ha $\varphi_1,\ldots,\varphi_n$ alaprendszer,
a $\Phi$ alapmátrixa az $x'=Ax$ homogén egyenletnek, akkor:
\begin{enumerate}[label=\arabic*)]
\item
\[
 \Phi'(t)=A(t)\Phi(t)
 \qquad\text{(mátrix-differenciálegyenlet)}.
\]
\item
\[
 \mathcal M_h=\left\{\sum_{k=1}^{n}c_k\varphi_k
 \;\middle|\;c_1,\ldots,c_n\in\mathbb R\right\},
\]
valamint, oszlopvektorral jelölve:
\[
 \mathcal M_h=\{\Phi\,c\mid c\in\mathbb R^n\}.
\]
\end{enumerate}
\end{matbox}

\subsection{Inhomogén lineáris rendszerek}
\textbf{Inhomogén egyenlet:}
\[
 x'=Ax+b,\qquad b\ne0.\tag{IH}
\]
\textbf{Észrevétel.} Ha $\psi_1,\psi_2$ az (IH) megoldásai,
akkor $\psi_1-\psi_2$ megoldása az $x'=Ax$ homogén egyenletnek.

\textbf{Tétel.} Legyen $\psi_p$ az inhomogén egyenlet egy
\textbf{partikuláris megoldása}. Ekkor az inhomogén egyenlet
általános megoldása:
\[
 \psi=\Phi\,c+\psi_p,
\]
ahol $\Phi$ a homogén egyenlet alapmátrixa és
$c\in\mathbb R^n$ oszlopvektor.

\subsubsection*{A partikuláris megoldás előállítása: az állandók variálása}
Keressük $\psi_p$-t a következő alakban:
\[
 \psi_p(t)=\Phi(t)c(t).
\]
Ekkor
\begin{align*}
 \psi_p'(t)
 &=\Phi'(t)c(t)+\Phi(t)c'(t)\\
 &=A(t)\Phi(t)c(t)+\Phi(t)c'(t)\\
 &=A(t)\psi_p(t)+b(t)\\
 &=A(t)\Phi(t)c(t)+b(t).
\end{align*}
A két $A(t)\Phi(t)c(t)$ tag kiesik. Az alapmátrix invertálható,
ezért $\Phi^{-1}(t)$ létezik, és
\[
 c'(t)=\Phi^{-1}(t)b(t)
 \quad\Longrightarrow\quad
 c(t)=\int_\tau^t\Phi^{-1}(s)b(s)\,ds.
\]

\begin{matbox}
\textbf{Tétel (Cauchy-formula).} Legyen $\Phi$ az $x'=Ax$
homogén egyenlet egy alapmátrixa. Minden
$(\tau,\xi)\in I\times\mathbb R^n$ esetén az
\[
 x'=Ax+b,\qquad x(\tau)=\xi
\]
kezdetiérték-probléma globálisan egyértelműen oldható meg,
és a $\psi$ teljes megoldás:
\begin{align*}
 \psi(t)&=\Phi(t)\Phi^{-1}(\tau)\xi
 +\Phi(t)\underbrace{\int_\tau^t
             \Phi^{-1}(s)b(s)\,ds}_{\psi_p(t)}
 \qquad(t\in I).
\end{align*}
\end{matbox}

\textbf{Megjegyzések.}
\begin{enumerate}[label=\arabic*)]
\item Az inhomogén kezdetiérték-probléma megoldásához
\textbf{csak} $\Phi$ ismerete kell (de $\Phi$ meghatározása nehéz feladat).
\item $\Phi$ explicit előállítása \textbf{csak (!)}
állandó együtthatós esetben lehetséges (lásd később).
\end{enumerate}

\subsection{Magasabbrendű, n-edrendű explicit differenciálegyenletek}
\textit{Hivatkozás: 08\_An\_alk\_2022\_osz.pdf.}

Legyen $n\in\mathbb N$, $h:\mathbb R\times\mathbb R^n\to\mathbb R$
folytonos függvény. Keressünk $I\subseteq\mathbb R$ nyílt intervallumon
$\varphi:I\to\mathbb R$, $\varphi\in C^n$ függvényt úgy, hogy
\[
 \bigl(t,\varphi(t),\varphi'(t),\ldots,
              \varphi^{(n-1)}(t)\bigr)\in D_h
 \qquad(t\in I),
\]
és
\[
 \varphi^{(n)}(t)
 =h\bigl(t,\varphi(t),\varphi'(t),\ldots,
                    \varphi^{(n-1)}(t)\bigr).
\]

\textbf{Röviden:}
\[
 y^{(n)}=h\circ\bigl(\operatorname{id},y,y',\ldots,
             y^{(n-1)}\bigr).\tag{\#}
\]
Ez \textbf{explicit}, mert az $n$-edik derivált explicit kifejezhető.

Legyen
\[
 (\tau,\xi)\in D_h\subseteq\mathbb R\times\mathbb R^n,
 \qquad\xi=(\xi_1,\ldots,\xi_n).
\]
Keressük az olyan $y$-t, amelyre
\[
 y(\tau)=\xi_1,\qquad y'(\tau)=\xi_2,
 \quad\ldots,\quad y^{(n-1)}(\tau)=\xi_n.
\]
Ez az \textbf{$n$-edrendű kezdetiérték-probléma}.

\textbf{Megjegyzés.} A mellékfeltételek más módon is megadhatók,
például peremfeltételekkel (sok eredmény született velük
kapcsolatban; lásd Tóth--Simon).

\subsection{Átviteli elv}
\textbf{Észrevétel.} Legyen $\varphi:I\to\mathbb R$,
$\varphi\in C^n$ az (\#) egyenlet megoldása. Legyen
\[
 \psi:=\begin{bmatrix}
     \varphi\\\varphi'\\\vdots\\\varphi^{(n-1)}
 \end{bmatrix}:I\to\mathbb R^n,
 \qquad
 \psi=\begin{bmatrix}\psi_1\\\vdots\\\psi_n\end{bmatrix}.
\]
Ekkor
\[
 \psi'=\begin{bmatrix}
 \varphi'\\\varphi''\\\vdots\\\varphi^{(n)}
 \end{bmatrix}
 =\begin{bmatrix}
 \psi_2\\\psi_3\\\vdots\\
 h\circ(\operatorname{id},\psi_1,\ldots,\psi_{n-1})
 \end{bmatrix}.
\]
% FORRÁSHŰ, 25. PDF-oldal: a kézirat az előző mátrixban
% a h argumentumát psi_{n-1}-nél zárja, és az f állapotvektora u_1-gyel indul.
% A matematikailag helyesnek vélt változatok az ellenorzes.md-ben szerepelnek.
Definiáljuk
\[
 f(t,u_1,\ldots,u_n):=
 \begin{bmatrix}
 u_1\\u_2\\\vdots\\h(t,u_1,\ldots,u_n)
 \end{bmatrix}\in\mathbb R^n.
\]
Ezért
\[
 \psi'=f\circ(\operatorname{id},\psi).
\]

\begin{matbox}
\textbf{Tétel (átviteli elv).}
\[
 \underbrace{y^{(n)}=h\circ(\operatorname{id},y,y',\ldots,
                         y^{(n-1)})}_{n\text{-edrendű d.e.}}
 \quad\Longleftrightarrow\quad
 \underbrace{x'=f\circ(\operatorname{id},x)}_{\text{elsőrendű d.e.r.}}.
\]
\end{matbox}

\begin{matbox}
\textbf{Tétel.} Tegyük fel, hogy $h\in C$, és $h$ a második
változójában lokális Lipschitz-feltételnek tesz eleget.
Ekkor minden $(\tau,\xi)\in D_h\subseteq\mathbb R\times\mathbb R^n$
pontban a hozzá tartozó kezdetiérték-probléma globálisan
egyértelműen oldható meg, és minden teljes megoldás
határtól határig halad.
\end{matbox}

\subsection{n-edrendű lineáris differenciálegyenletek}
\textbf{Általános alak} ($n\in\mathbb N$):
\[
 y^{(n)}+a_{n-1}y^{(n-1)}+\cdots+a_1y'+a_0y=b.
\]
Itt $a_i,b:I\to\mathbb R$ folytonos függvények,
$I\subseteq\mathbb R$ nyílt intervallum.
\textit{(Az eredeti sorban az indexhatár $i=0,\ldots,n$;
a képlet $a_0,\ldots,a_{n-1}$ együtthatókat tartalmaz.
Az eltérést az ellenőrzési jelentés tárgyalja.)}

\begin{itemize}
\item $b\equiv0$: \textbf{homogén egyenlet}.
\item $b\ne0$: \textbf{inhomogén egyenlet}.
\end{itemize}

\begin{matbox}
\textbf{Tétel.}
\begin{enumerate}[label=\arabic*)]
\item Minden $(\tau,\xi)\in I\times\mathbb R^n$ esetén
az inhomogén differenciálegyenletre vonatkozó kezdetiérték-probléma
globálisan egyértelműen oldható meg, és minden teljes megoldás
értelmezési tartománya az egész $I$.
\item A homogén egyenletnek van $n$ darab lineárisan független megoldása,
és minden megoldás ezek lineáris kombinációja.
\item $\varphi_1,\varphi_2,\ldots,\varphi_n$ homogén megoldások
lineárisan függetlenek akkor és csak akkor, ha
\[
 \exists t_0\in I:\qquad\det\bigl(W(t_0)\bigr)\ne0,
\]
ahol a \textbf{Wronski-féle mátrix}:
\[
 W(t)=\begin{bmatrix}
 \varphi_1(t)&\cdots&\varphi_n(t)\\
 \varphi_1'(t)&\cdots&\varphi_n'(t)\\
 \vdots&&\vdots\\
 \varphi_1^{(n-1)}(t)&\cdots&\varphi_n^{(n-1)}(t)
 \end{bmatrix}.
\]
A $\det W(t)$ a \textbf{Wronski-determináns}.
\textit{Bizonyítás: lásd Tóth--Simon.}
\end{enumerate}
\end{matbox}

\begin{matbox}
\textbf{Tétel (szuperpozíció).}
Ha $\psi_p$ az inhomogén egyenlet partikuláris megoldása,
akkor az inhomogén egyenlet általános megoldása
\[
 \psi=\varphi+\psi_p,
\]
ahol $\varphi$ a homogén egyenlet általános megoldása.
\end{matbox}

\subsubsection*{Állandók variálása az $n$-edrendű lineáris egyenletnél}
\begin{matbox}
\textbf{Tétel (állandók variálása).} Tekintsük
\[
 y^{(n)}+a_{n-1}y^{(n-1)}+\cdots+a_1y'+a_0y=b,
\]
ahol az $a_i,b:I\to\mathbb R$ függvények folytonosak
($i=0,\ldots,n-1$).
Legyen $\varphi_1,\ldots,\varphi_n$ a homogén egyenlet
\textbf{alaprendszere}. Ekkor az inhomogén egyenlet egy
$\psi_p$ partikuláris megoldása:
\[
 \psi_p(t)=c_1(t)\varphi_1(t)+\cdots+c_n(t)\varphi_n(t),
 \qquad t\in I,
\]
ahol
\[
 c'(t)=\begin{bmatrix}
 c_1'(t)\\\vdots\\c_n'(t)
 \end{bmatrix},
 \qquad\Phi(t)=W(t),\qquad
 W(t)c'(t)=\begin{bmatrix}
 0\\0\\\vdots\\b(t)
 \end{bmatrix}.
\]
\end{matbox}

\textbf{Megjegyzés.} Az inhomogén megoldáshoz elég a homogén
egyenlet egy alaprendszerét ismerni. Az explicit előállítás
\textbf{csak} állandó együtthatós esetben van általános módszer.
